\centerline{\bf \Huge Счётный планик II. КР}
\centerline{\bf \Huge }

\begin{flushright}
        \scriptsize{\textbf{yy}}
\end{flushright}


\problem{\textbf{1}} В трапецию $ABCD$ с основаниями $AD$ и $BC$ вписана окружность с центром в $O$.

а) Докажите, что $\sin \angle AOD = \sin \angle BOC$.

б) Найдите площадь трапеции, если $\angle BAD = 90^{\circ}$, а основания равны 5 и 7.

\problem{\textbf{2}} Окружность, вписанная в трапецию $ABCD$, касается её боковых сторон $AB$ и $CD$ в точках
$M$ и $N$ соответственно. Известно, что $AM = 6MB$ и $2DN = 3CN$.

а) Докажите, что $AD = 3BC$.

б) Найдите длину отрезка $MN$, если радиус окружности равен $\sqrt{105}$

\problem{\textbf{3}} Дана равнобедренная трапеция $ABCD$ с основаниями $AD$ и $BC$. Биссектрисы углов $BAD$
и $BCD$ пересекаются в точке $O$. Точки $M$ и $N$ отмечены на боковых сторонах $AB$ и $CD$
соответственно. Известно, что $AM = MO, CN = NO$.

а) Докажите, что точки $M$, $N$ и $O$ лежат на одной прямой.

б) Найдите $AM : MB$, если известно, что $AO = OC$ и $BC : AD = 1 : 7$